%% 第 12 章：Jordan 链
%% 本文件由 tools/md2tex.py 自动生成，请勿直接编辑。

\chapter{Jordan 链}


\section{从广义特征空间得到幂零算子}

固定一个特征值 $\lambda$，并限制到其广义特征空间：

\[
G_\lambda.
\]

令：

\[
N=(T-\lambda I)|_{G_\lambda}.
\]

则 $N$ 是 $G_\lambda$ 上的幂零算子。

因此，Jordan 标准型的问题可以转化为：

\begin{jinsight}
如何描述一个幂零算子的内部结构？
\end{jinsight}

\begin{jdefn}{Jordan 链}{r:13:1}
长度为 $r$ 的 Jordan 链是一列向量：

\[
(v_1,\ldots,v_r),
\]

满足：

\[
Nv_1=0,
\]

以及：

\[
Nv_j=v_{j-1},
\qquad
j=2,\ldots,r.
\]

等价地：

\[
Tv_1=\lambda v_1,
\]

\[
Tv_j=\lambda v_j+v_{j-1},
\qquad
j=2,\ldots,r.
\]

其中 $v_1$ 是普通特征向量；后续向量是广义特征向量，但会逐层经 $N$ 回落到前一个向量。
\end{jdefn}

\begin{jprop}{单条 Jordan 链中的向量线性无关}{r:13:2}
若存在向量 $v$，满足：

\[
N^{r}v=0,
\qquad
N^{r-1}v\ne0,
\]

则：

\[
v,Nv,\ldots,N^{r-1}v
\]

线性无关。
\end{jprop}

\begin{jproof}{证明思路}
假设这列向量存在非平凡线性关系。

取其中最低次、且系数不为零的项。对整条关系施加恰当次数的 $N$，所有更高次项都会因 $N^rv=0$ 而消失，只留下该系数乘以：

\[
N^{r-1}v.
\]

但：

\[
N^{r-1}v\ne0.
\]

这与该系数非零矛盾。

因此，一条由链头不断施加 $N$ 生成的完整链，内部一定线性无关。
\end{jproof}

\section{Jordan 链基的构造原则}

一般的幂零算子不一定只需要一条链，也就是说，这个幂零算子的特征值对应的广义特征空间中，可能存在多条完整的链。

不能把：

\[
\ker N,\ker N^2,\ldots
\]

中的所有向量直接全部取来做基，因为这些空间彼此嵌套，其中包含重复方向。

正确做法是：从零空间链的高层开始，挑选真正新增的独立方向作为链头；再沿 $N$ 向下生成链。

设：

\[
K_k=\ker N^k,
\qquad
K_0=\{0\}.
\]

则：

\[
\dim K_k-\dim K_{k-1}
\]

等于长度至少为 $k$ 的 Jordan 链数量。

而：

\[
\bigl(\dim K_k-\dim K_{k-1}\bigr)
-
\bigl(\dim K_{k+1}-\dim K_k\bigr)
\]

等于长度恰好为 $k$ 的 Jordan 链数量。

\begin{jthm}{Jordan 链基存在定理}{r:13:3}
每个有限维幂零算子都存在一组由 Jordan 链组成的基。

在 Jordan 标准型的语境中，先固定一个特征值 $\lambda$，并令：

\[
W=G_\lambda,
\]

\[
N=(T-\lambda I)|_W:W\to W.
\]

此时 $N$ 是 $W$ 上的幂零算子。以下构造的目标是为 $W$ 找到一组 Jordan 链基，而不是直接为原始空间 $V$ 找基。
\end{jthm}

\begin{jidea}{构造的目标}
我们希望选出若干条链：

\[
c,\ Nc,\ N^2c,\ldots
\]

使每条链内部线性无关，并且所有链合在一起后，既不重复，也恰好张成整个 $W$。

这里将 $c$ 称为一条链的生成端：不断对它施加 $N$，便得到整条链。
\end{jidea}

\section{第一步：确定最长链的长度}

取最小的正整数 $r$，使：

\[
N^r=0.
\]

这意味着：对 $W$ 中任何向量连续作用 $N$ 共 $r$ 次，结果都会变为零；但至少存在某个向量，需要恰好 $r$ 次才变成零。

记：

\[
K_k=\ker N^k.
\]

此时：

\[
K_r=W.
\]

因为 $N^r=0$ 表示 $W$ 中每个向量都属于 $\ker N^r$。

\subsection{第二步：先构造最长的链}

$K_{r-1}$ 包含所有经过至多 $r-1$ 次作用就会变为零的向量。

先取 $K_{r-1}$ 的一组基，再将它扩充为 $W$ 的一组基。新补入的向量张成某个子空间，记为 $C_r$：

\[
W=K_{r-1}\oplus C_r.
\]

任取非零 $c\in C_r$。因为：

\[
c\in W=K_r,
\qquad
c\notin K_{r-1},
\]

所以：

\[
N^rc=0,
\qquad
N^{r-1}c\ne0.
\]

因此，$c$ 生成一条长度恰为 $r$ 的链：

\[
c,\ Nc,\ N^2c,\ldots,N^{r-1}c.
\]

对 $C_r$ 的每个基向量都这样做，便得到全部最长的 Jordan 链。

\subsection{第三步：构造较短的链}

接下来处理长度为 $r-1$ 的链。

不能直接从 $K_{r-1}$ 中随便挑向量，因为其中已经包含刚才长链向下作用得到的向量。例如，若：

\[
c,\ Nc,\ldots,N^{r-1}c
\]

是一条长链，那么 $Nc$ 已经位于 $K_{r-1}$ 中；不能再把 $Nc$ 当作另一条新链的生成端。

因此，在第 $k$ 层寻找长度为 $k$ 的链时，要先排除两类已经解释过的方向：

\begin{jitem}
\item %
$K_{k-1}$ 中的向量：它们在不足 $k$ 次作用后已经归零，长度不够；
\item %
已有较长链向下作用后得到的向量：它们已经属于旧链的一部分，不应重复使用。
\end{jitem}

然后，在 $K_k$ 中选取一批无法由这些旧方向表示的新向量。

具体做法仍是熟悉的“扩充基”：

\begin{jenum}
\item %
先取全部旧方向的一组基；
\item %
将它扩充为 $K_k$ 的一组基；
\item %
新补入的向量，就是长度恰为 $k$ 的新链生成端。
\end{jenum}

对每一个新生成端 $c$，依次写出：

\[
c,\ Nc,\ldots,N^{k-1}c.
\]

由于：

\[
c\notin K_{k-1},
\]

它不会提前归零，因此这确实是一条长度为 $k$ 的完整链。

按：

\[
r,r-1,\ldots,1
\]

的顺序重复这个过程，直到所有层都处理完。

\subsection{为什么不同链之间不会线性相关}

单条链内部的线性无关性已经证明过。

不同链之间不会产生冗余，原因在于：每次选择新的生成端时，都先排除了已经属于旧链或能够由旧链表示的方向；随后再通过扩充基，选取真正独立的新方向。

因此，新链不会重复旧链，也不能由已经构造出的链线性表示。

最终，所有链共同张成：

\[
K_1,\ K_2,\ \ldots,\ K_r=W.
\]

也就是说，它们共同张成整个广义特征空间 $W$。

同时，构造过程中每增加一条链，增加的都是此前未出现的独立方向；所有链向量的总数最终恰好等于：

\[
\dim W.
\]

因此，这些向量张成 $W$，且数量恰好等于 $W$ 的维数。它们合起来便构成 $W$ 的一组基。

\subsection{回到原始算子}

对每个特征值 $\lambda$，都在其广义特征空间：

\[
G_\lambda
\]

上重复上述构造。

最后，由主分解：

\[
V=\bigoplus_\lambda G_\lambda,
\]

将所有广义特征空间中的 Jordan 链基合并，便得到原始空间 $V$ 的一组 Jordan 基。

\begin{jinsight}
Jordan 链基的构造不是任意挑选若干条链，而是在每个广义特征空间中，从最长链开始，逐层排除旧方向，再补入新的独立生成端。
\end{jinsight}
